% ============================================================
%  CUADERNILLO DE ESTUDIO — CLASE 6
%  Seguridad de la Información y de Sistemas
%  Lic. en Gestión de Negocios Digitales — UCA
%  Compilar con: pdflatex (dos pasadas para el índice)
% ============================================================

\documentclass[12pt,a4paper]{article}

% ---------- Paquetes ----------
\usepackage[utf8]{inputenc}
\usepackage[spanish]{babel}
\usepackage[T1]{fontenc}
\usepackage{lmodern}
\usepackage[margin=2.2cm,headheight=14pt]{geometry}
\usepackage{amsmath}
\usepackage{amssymb}
\usepackage{enumitem}
\usepackage{booktabs}
\usepackage{array}
\usepackage{tabularx}
\usepackage[table]{xcolor} % Necesario para \cellcolor en tablas
\usepackage{tcolorbox}
\usepackage{graphicx}
\usepackage{fancyhdr}
\usepackage{titlesec}
\usepackage{hyperref}
\usepackage{multicol}
\usepackage{pifont}
\usepackage{wasysym}
\usepackage{tikz}
\usetikzlibrary{positioning,arrows.meta,shapes.geometric}

% ---------- Colores ----------
\definecolor{azulUCA}{HTML}{003366}
\definecolor{doradoUCA}{HTML}{B8860B}
\definecolor{grisClaro}{HTML}{F2F2F2}
\definecolor{rojoAlerta}{HTML}{C0392B}
\definecolor{verdeOK}{HTML}{27AE60}
\definecolor{azulCielo}{HTML}{2980B9}
\definecolor{naranjaAmenaza}{HTML}{D35400}
\definecolor{violetaParcial}{HTML}{6C3483}

% Colores para la Matriz de Riesgos (Mapa de Calor)
\definecolor{riskLow}{HTML}{A9DFBF}
\definecolor{riskMed}{HTML}{F9E79F}
\definecolor{riskHigh}{HTML}{F5CBA7}
\definecolor{riskCrit}{HTML}{F1948A}

% ---------- tcolorbox estilos ----------
\tcbuselibrary{skins,breakable}

\newtcolorbox{cajaTitulo}[1][]{
  colback=azulUCA, colframe=azulUCA, coltext=white,
  fonttitle=\Large\bfseries, arc=4mm, #1
}

\newtcolorbox{cajaDefinicion}{
  colback=azulCielo!10, colframe=azulCielo,
  fonttitle=\bfseries, title=Definici\'on, arc=3mm, breakable
}

\newtcolorbox{cajaConcepto}{
  colback=doradoUCA!10, colframe=doradoUCA,
  fonttitle=\bfseries, title=Concepto clave, arc=3mm, breakable
}

\newtcolorbox{cajaActividad}{
  colback=grisClaro, colframe=azulUCA,
  fonttitle=\bfseries\color{azulUCA},
  title={\ding{45}\; Actividad}, arc=3mm, breakable
}

\newtcolorbox{cajaParcial}{
  colback=violetaParcial!8, colframe=violetaParcial,
  fonttitle=\bfseries\color{violetaParcial},
  title={\ding{115}\; PARCIAL 1}, arc=3mm, breakable
}

\newtcolorbox{cajaNota}{
  colback=rojoAlerta!8, colframe=rojoAlerta,
  fonttitle=\bfseries\color{rojoAlerta}, title=Nota de Negocio, arc=3mm, breakable
}

% ---------- Header / Footer ----------
\pagestyle{fancy}
\fancyhf{}
\fancyhead[L]{\small\textcolor{azulUCA}{Seguridad de la Informaci\'on y de Sistemas}}
\fancyhead[R]{\small\textcolor{azulUCA}{Cuadernillo --- Clase 6}}
\fancyfoot[C]{\thepage}
\fancyfoot[L]{\small\textcolor{gray}{UCA}}
\renewcommand{\headrulewidth}{0.4pt}
\renewcommand{\footrulewidth}{0.4pt}

% ---------- Títulos de sección ----------
\titleformat{\section}
  {\color{azulUCA}\Large\bfseries}
  {\thesection.}{0.6em}{}
\titleformat{\subsection}
  {\color{azulCielo}\large\bfseries}
  {\thesubsection}{0.6em}{}

% ---------- Enlaces ----------
\hypersetup{colorlinks=true, linkcolor=azulUCA, urlcolor=azulCielo}

% ---------- Checkbox ----------
\newcommand{\checkboxVacio}{$\square$\;}
\newcommand{\checkMark}{\ding{51}}

% ============================================================
\begin{document}

% ============================================================
%  PORTADA
% ============================================================
\begin{titlepage}
\centering
\vspace*{1cm}

\begin{cajaTitulo}[width=\textwidth, halign=center]
  SEGURIDAD DE LA INFORMACI\'ON Y DE SISTEMAS\\[4pt]
  {\large Licenciatura en Gesti\'on de Negocios Digitales --- UCA}
\end{cajaTitulo}

\vspace{1.2cm}

{\Huge\bfseries\color{azulUCA} CUADERNILLO DE ESTUDIO}\\[8pt]
{\LARGE\color{doradoUCA} CLASE 6}\\[10pt]
{\Large\itshape ``Gesti\'on de riesgos II:\\[2pt] BIA, controles e ISO 27001''}\\[6pt]
{\large\color{violetaParcial}\ding{115}\; Incluye preparaci\'on para el PARCIAL 1}\\[1.5cm]

\begin{tabular}{rl}
\textbf{Profesor:}   & Mg.\ Ing.\ Rodrigo Atilio Elgueta \\[4pt]
\textbf{Unidad:}     & Cierre de Unidad 2 \\[4pt]
\end{tabular}

\vfill

\begin{tcolorbox}[colback=grisClaro, colframe=azulUCA, arc=3mm, width=0.9\textwidth]
  \centering
  {\bfseries Datos del/la estudiante}\\[10pt]
  Nombre y apellido: \underline{\hspace{5cm}}\\[10pt]
  Grupo N.\textdegree: \underline{\hspace{5cm}}\\[10pt]
  Negocio Digital (TP): \underline{\hspace{5cm}}\\[10pt]
  Fecha: \underline{\hspace{5cm}}
\end{tcolorbox}

\vspace{1cm}
{\small\textcolor{gray}{Material de uso exclusivo para la cursada.}}
\end{titlepage}

% ============================================================
%  ÍNDICE
% ============================================================
\tableofcontents
\newpage

% ============================================================
%  SECCIÓN 1 — REPASO EXPRESS
% ============================================================
\section{Repaso express: lo que traemos de la Clase 5}

Antes de avanzar, aseguremos los cimientos. En la Clase 5 vimos el \textbf{ciclo de gesti\'on de riesgos}:

\begin{center}
\begin{tikzpicture}[
  node distance=0.6cm,
  box/.style={rectangle, draw=azulUCA, fill=azulUCA!10,
              text width=3cm, align=center, rounded corners=3pt,
              font=\small\bfseries},
  arr/.style={-{Stealth[length=3mm]}, thick, azulUCA}
]
  \node[box] (a) {1. IDENTIFICAR\\{\scriptsize ¿Qu\'e puede salir mal?}};
  \node[box, right=1cm of a] (b) {2. ANALIZAR\\{\scriptsize Probabilidad e Impacto}};
  \node[box, right=1cm of b] (c) {3. EVALUAR\\{\scriptsize Matriz P$\times$I}};
  \node[box, right=1cm of c] (d) {4. TRATAR\\{\scriptsize Mitigar/Transferir/Aceptar/Evitar}};
  \draw[arr] (a) -- (b);
  \draw[arr] (b) -- (c);
  \draw[arr] (c) -- (d);
\end{tikzpicture}
\end{center}

Hoy profundizamos en dos piezas clave de ese ciclo:
\begin{itemize}[leftmargin=2em]
    \item El \textbf{an\'alisis de impacto} (BIA), que le pone n\'umeros al da\~no.
    \item Los \textbf{controles}, que son las herramientas concretas para \textbf{tratar} los riesgos.
\end{itemize}

\begin{cajaConcepto}[title={Pregunta disparadora}]
Si el servidor de facturaci\'on de tu negocio se cae 4 horas un s\'abado a la noche, ¿cu\'anto dinero perd\'es? ¿Y si se cae un martes a las 11 de la ma\~nana? \textbf{No todos los procesos valen lo mismo, ni todas las ca\'idas duelen igual.} Eso es lo que mide el BIA.
\end{cajaConcepto}

\newpage
% ============================================================
%  SECCIÓN 2 — BIA
% ============================================================
\section{BIA: An\'alisis de Impacto al Negocio}

\begin{cajaDefinicion}
El \textbf{Business Impact Analysis (BIA)} es el proceso de identificar los procesos cr\'iticos del negocio y cuantificar el impacto que tendr\'ia su interrupci\'on, con el objetivo de definir prioridades de recuperaci\'on y continuidad.
\end{cajaDefinicion}

\subsection{¿Para qu\'e sirve?}
\begin{itemize}[leftmargin=2em]
    \item \textbf{Priorizar:} no todo se recupera al mismo tiempo; primero lo cr\'itico.
    \item \textbf{Dimensionar inversiones:} justificar cu\'anto gastar en backup, redundancia o contingencia.
    \item \textbf{Definir objetivos de recuperaci\'on:} cu\'anto tiempo podemos tolerar sin cada proceso.
\end{itemize}

\subsection{Los tres conceptos fundamentales del BIA}

\renewcommand{\arraystretch}{1.6}
\begin{tabularx}{\textwidth}{|>{\bfseries\color{azulUCA}}p{2.6cm}|X|p{4.5cm}|}
\hline
\rowcolor{grisClaro}
\textbf{Sigla} & \textbf{¿Qu\'e pregunta?} & \textbf{Ejemplo} \\
\hline
\textbf{MTD} & ¿Cu\'anto tiempo m\'aximo puede estar ca\'ido este proceso antes de que el da\~no sea inaceptable? & ``La pasarela de pago no puede estar ca\'ida m\'as de 2 horas.'' \\
\hline
\textbf{RTO} & ¿En cu\'anto tiempo nos \textbf{proponemos} recuperarlo tras un incidente? & ``Objetivo: volver a operar en 1 hora.'' \\
\hline
\textbf{RPO} & ¿Cu\'antos \textbf{datos} podemos permitirnos perder? (cada cu\'anto hago backup) & ``Backup cada 15 min $\rightarrow$ pierdo m\'ax. 15 min de datos.'' \\
\hline
\end{tabularx}

\begin{cajaNota}
\textbf{Regla pr\'actica:} Siempre debe cumplirse que \textbf{RTO $\leq$ MTD}. Si tu proceso tolera m\'aximo 2 horas ca\'ido (MTD) pero tu plan de recuperaci\'on tarda 6 horas (RTO), ten\'es un problema de dise\~no. El RPO, adem\'as, define la frecuencia m\'inima de tus backups.
\end{cajaNota}

\subsection{Plantilla de BIA simplificado}

Esta es la tabla que van a usar en el Parcial 1. Completala para cada proceso cr\'itico:

\renewcommand{\arraystretch}{1.8}
\begin{tabularx}{\textwidth}{|X|}
\hline
\textbf{Proceso / Activo:} \hrulefill \\[4pt]
\textbf{Impacto por hora de ca\'ida:} \hrulefill \quad \textbf{(Econ\'omico / Legal / Reputacional / Operativo)} \\[4pt]
\textbf{MTD (M\'ax. tiempo tolerable):} \hrulefill \\[4pt]
\textbf{RTO (Objetivo de recuperaci\'on):} \hrulefill \\[4pt]
\textbf{RPO (P\'erdida m\'ax. de datos aceptable):} \hrulefill \\[4pt]
\textbf{Prioridad (Cr\'itico / Alto / Medio / Bajo):} \hrulefill \\[4pt]
\hline
\end{tabularx}

\begin{cajaActividad}[title={Ejemplo guiado: e-commerce de indumentaria}]
\textbf{Proceso:} Pasarela de pago (checkout).\\
\textbf{Impacto por hora de ca\'ida:} Econ\'omico alto (se pierden ventas directas) + reputacional (clientes frustrados en redes).\\
\textbf{MTD:} 2 horas (en temporada alta, 30 minutos).\\
\textbf{RTO:} 1 hora.\\
\textbf{RPO:} 5 minutos (las transacciones no se pueden perder).\\
\textbf{Prioridad:} \textbf{CR\'ITICO}.
\end{cajaActividad}

\newpage
% ============================================================
%  SECCIÓN 3 — TIPOS DE CONTROLES
% ============================================================
\section{Tipos de controles y costo-beneficio}

Un \textbf{control} es cualquier medida (t\'ecnica, humana, f\'isica u organizativa) dise\~nada para reducir la probabilidad o el impacto de un riesgo.

\subsection{Las 4 categor\'ias de controles}

\renewcommand{\arraystretch}{1.5}
\begin{tabularx}{\textwidth}{|>{\bfseries\color{azulUCA}}p{3cm}|X|X|}
\hline
\rowcolor{grisClaro}
\textbf{Categor\'ia} & \textbf{¿Qu\'e abarca?} & \textbf{Ejemplos} \\
\hline
\textbf{T\'ecnicos} & Soluciones tecnol\'ogicas que previenen, detectan o corrigen. & Firewall, MFA, cifrado, antivirus, backups autom\'aticos, IDS/IPS. \\
\hline
\textbf{Humanos / Organizativos} & Acciones sobre las personas y los procesos. & Capacitaci\'on, concientizaci\'on, pol\'iticas, roles, procedimientos. \\
\hline
\textbf{F\'isicos} & Protecci\'on del espacio y el hardware. & C\'amaras, control de acceso al data center, cerco perimetral, UPS. \\
\hline
\textbf{Legales / Normativos} & Cumplimiento de leyes y contratos. & Pol\'itica de privacidad, NDAs, seguros, cl\'ausulas con proveedores. \\
\hline
\end{tabularx}

\subsection{Adem\'as: ¿previenen, detectan o corrigen?}
Cada control puede clasificarse seg\'un su \textbf{funci\'on en el tiempo}:
\begin{itemize}[leftmargin=2em]
    \item \textbf{Preventivo:} evita que el incidente ocurra (ej. MFA, capacitaci\'on).
    \item \textbf{Detectivo:} alerta cuando el incidente est\'a ocurriendo (ej. antivirus, c\'amaras, monitoreo).
    \item \textbf{Correctivo:} restaura la normalidad despu\'es del incidente (ej. backups, plan de respuesta).
\end{itemize}

\subsection{Criterio costo-beneficio}
\begin{cajaConcepto}[title={La pregunta de negocio}]
Un control solo se implementa si \textbf{el costo del control es menor que el impacto esperado del riesgo}.\\[4pt]
\textbf{F\'ormula conceptual:}
\begin{center}
$\text{Beneficio del control} \;\approx\; (\text{Impacto del riesgo}) - (\text{Costo del control})$
\end{center}
Si proteger un activo cuesta m\'as que el activo mismo, la decisi\'on correcta suele ser \textbf{aceptar} o \textbf{transferir} el riesgo, no mitigarlo.
\end{cajaConcepto}

\begin{cajaActividad}[title={Ejercicio r\'apido: justific\'a el costo-beneficio}]
\renewcommand{\arraystretch}{1.6}
\begin{tabularx}{\textwidth}{|X|}
\hline
\textbf{Riesgo:} Filtraci\'on de base de datos de 10.000 clientes por un bucket mal configurado. \\
\textbf{Impacto estimado:} Multa + p\'erdida de clientes $\approx$ \$500.000. \\
\textbf{Control propuesto:} Revisi\'on de configuraciones cloud + cifrado. Costo: \$30.000/a\~no. \\[4pt]
\textbf{¿Lo implement\'as? Justific\'a en 2-3 líneas:} \\[2.5cm]
\hline
\end{tabularx}
\end{cajaActividad}

\newpage
% ============================================================
%  SECCIÓN 4 — ISO 27001 / 27002
% ============================================================
\section{Introducci\'on a ISO/IEC 27001 y 27002}

Cuando una organizaci\'on quiere ordenar toda su seguridad de forma sistem\'atica, recurre a est\'andares internacionales. Los dos m\'as importantes son la familia ISO/IEC 27000.

\subsection{¿Qu\'e es cada una?}

\renewcommand{\arraystretch}{1.6}
\begin{tabularx}{\textwidth}{|>{\bfseries\color{azulUCA}}p{3cm}|X|X|}
\hline
\rowcolor{grisClaro}
 & \textbf{ISO/IEC 27001} & \textbf{ISO/IEC 27002} \\
\hline
\textbf{¿Qu\'e es?} & Norma de \textbf{requisitos} para un Sistema de Gesti\'on de Seguridad de la Informaci\'on (SGSI). & \textbf{Cat\'alogo de buenas pr\'acticas} y controles para implementar la seguridad. \\
\hline
\textbf{Rol} & Dice \textbf{QU\'E} hay que hacer (los requisitos del sistema de gesti\'on). & Sugiere \textbf{C\'OMO} hacerlo (lista de controles concretos). \\
\hline
\textbf{¿Certificable?} & \textbf{S\'i.} Una empresa puede certificarse en 27001. & \textbf{No} por sí sola; acompa\~na a la 27001. \\
\hline
\textbf{Analog\'ia} & Es como la ``constituci\'on'' que define que debe existir un sistema de seguridad. & Es como el ``manual de procedimientos'' con controles puntuales. \\
\hline
\end{tabularx}

\begin{cajaConcepto}[title={La relaci\'on entre ambas}]
La \textbf{27001} establece que la organizaci\'on debe gestionar sus riesgos de seguridad de forma sistem\'atica (planificar, implementar, verificar, mejorar). La \textbf{27002} le da el cat\'alogo de controles posibles (control de acceso, criptograf\'ia, seguridad f\'isica, gesti\'on de incidentes, etc.) para que elija cu\'ales aplicar seg\'un su an\'alisis de riesgos.
\end{cajaConcepto}

\subsection{¿Por qu\'e le importa a un negocio digital?}
\begin{itemize}[leftmargin=2em]
    \item \textbf{Confianza comercial:} estar certificado (o alineado) a ISO 27001 es una se\~nal fuerte para clientes y partners.
    \item \textbf{Orden interno:} da un marco para no ``apagar incendios'' sino gestionar la seguridad con m\'etodo.
    \item \textbf{Ventaja competitiva:} muchas licitaciones y contratos B2B exigen cumplimiento de est\'andares de seguridad.
\end{itemize}

\begin{cajaNota}
\textbf{Para tu TP Integrador:} no necesit\'as certificar ISO 27001 (eso es costoso y lleva meses). Pero \textbf{alinear} tu programa de seguridad con la l\'ogica de la norma (identificar riesgos $\rightarrow$ elegir controles $\rightarrow$ documentar pol\'iticas $\rightarrow$ mejorar continuamente) le da rigor profesional a tu trabajo.
\end{cajaNota}

\newpage
% ============================================================
%  SECCIÓN 5 — CIERRE DE UNIDAD 2
% ============================================================
\section{Cierre de Unidad 2: mapa integrador}

Con esta clase cerramos la \textbf{Unidad 2: Gesti\'on de Riesgos}. As\'i se conecta todo lo visto entre la Unidad 1 y la Unidad 2:

\begin{center}
\begin{tikzpicture}[
  node distance=1.1cm and 1.8cm,
  box/.style={rectangle, draw=azulUCA, fill=azulUCA!10,
              text width=3.4cm, align=center, rounded corners=3pt,
              font=\small\bfseries},
  boxsmall/.style={rectangle, draw=azulCielo, fill=azulCielo!10,
              text width=3cm, align=center, rounded corners=3pt,
              font=\scriptsize},
  arr/.style={-{Stealth[length=2.5mm]}, thick, azulUCA}
]
  % Fila 1: Unidad 1
  \node[box] (cid) {UNIDAD 1\\Pilares CID y\\5 conceptos};
  \node[box, right=2.2cm of cid] (amen) {UNIDAD 1\\Amenazas y\\herramientas};
  \node[box, right=2.2cm of amen] (cripto) {UNIDAD 1\\Criptograf\'ia\\y confianza};

  % Fila 2: Unidad 2
  \node[boxsmall, below=1.4cm of cid] (matriz) {UNIDAD 2\\Matriz P$\times$I\\y ciclo de riesgos};
  \node[boxsmall, below=1.4cm of amen] (bia) {UNIDAD 2\\BIA y\\controles};
  \node[boxsmall, below=1.4cm of cripto] (iso) {UNIDAD 2\\ISO 27001/27002\\(marco de gesti\'on)};

  \draw[arr] (cid.south) -- (matriz.north);
  \draw[arr] (amen.south) -- (bia.north);
  \draw[arr] (cripto.south) -- (iso.north);
  \draw[arr, dashed] (matriz.east) -- (bia.west);
  \draw[arr, dashed] (bia.east) -- (iso.west);
\end{tikzpicture}
\end{center}

\begin{cajaConcepto}[title={Lo que nos llevamos de la Unidad 2}]
\begin{itemize}[leftmargin=2em]
    \item Aplicar el \textbf{ciclo de gesti\'on de riesgos} (identificar, analizar, evaluar, tratar).
    \item Construir y leer una \textbf{matriz de probabilidad e impacto}.
    \item Cuantificar el da\~no con un \textbf{BIA} (MTD, RTO, RPO).
    \item Elegir \textbf{controles} con criterio \textbf{costo-beneficio}.
    \item Conocer el marco \textbf{ISO/IEC 27001/27002} como referencia profesional.
\end{itemize}
\end{cajaConcepto}

\newpage
% ============================================================
%  SECCIÓN 6 — PARCIAL 1: FORMATO Y RÚBRICA
% ============================================================
\section{PARCIAL 1: formato, r\'ubrica y recomendaciones}

\begin{cajaParcial}[title={Ficha t\'ecnica del Parcial 1}]
\renewcommand{\arraystretch}{1.5}
\begin{tabularx}{\textwidth}{>{\bfseries}p{4cm}X}
\textbf{Modalidad:} & Grupal (mismos grupos del TP Integrador), en clase. \\
\textbf{Duraci\'on:} & 70 minutos. \\
\textbf{Cobertura:} & Unidades 1 y 2. \\
\textbf{Caso:} & Un negocio digital \textbf{nuevo}, no visto antes en clase. \\
\textbf{Nota:} & 1 a 10. Se aprueba con 4. \\
\end{tabularx}
\end{cajaParcial}

\subsection{¿Qu\'e deben entregar?}
Cada grupo recibe un caso nuevo y debe producir:
\begin{enumerate}[leftmargin=2em]
    \item Una \textbf{matriz de riesgos} (identificar activos, amenazas, vulnerabilidades y ubicarlos en la matriz P$\times$I).
    \item Un \textbf{BIA simplificado} de los procesos cr\'iticos del caso.
    \item \textbf{3 controles priorizados}, cada uno con su \textbf{justificaci\'on costo-beneficio}.
\end{enumerate}

\subsection{R\'ubrica de evaluaci\'on (sobre 10 puntos)}

\renewcommand{\arraystretch}{1.5}
\begin{tabularx}{\textwidth}{|X|>{\centering\arraybackslash}p{3cm}|}
\hline
\rowcolor{violetaParcial!20}
\textbf{¿Qu\'e se eval\'ua?} & \textbf{Puntos} \\
\hline
Identificaci\'on de activos y riesgos del caso & hasta 3 \\
\hline
Correcta aplicaci\'on de la matriz de probabilidad e impacto & hasta 3 \\
\hline
Pertinencia y justificaci\'on costo-beneficio de los controles propuestos & hasta 2 \\
\hline
Comprensi\'on individual demostrada en las preguntas cruzadas de cierre & hasta 2 \\
\hline
\rowcolor{grisClaro}
\textbf{TOTAL} & \textbf{10} \\
\hline
\end{tabularx}

\subsection{Las preguntas cruzadas individuales}
Al final, el docente har\'a preguntas individuales a cada integrante para verificar la comprensi\'on \textbf{personal} dentro del trabajo grupal. Ejemplos de lo que pueden preguntarles:
\begin{itemize}[leftmargin=2em]
    \item ¿Por qu\'e ubicaste ese riesgo en ``Alto'' y no en ``Medio''?
    \item ¿Qu\'e diferencia hay entre el control que propusiste y uno preventivo?
    \item ¿C\'omo justific\'as el costo-beneficio de ese control?
    \item Si el BIA dice que el proceso tolera 4 horas ca\'ido, ¿por qu\'e propusiste un RTO de 8?
\end{itemize}

\begin{cajaNota}[title={Consejos para el d\'ia del Parcial}]
\begin{enumerate}[leftmargin=2em]
    \item \textbf{Distribuyan roles r\'apido:} uno arma la matriz, otro el BIA, otro los controles, todos revisan.
    \item \textbf{Lean el caso dos veces} antes de empezar; la mayor\'ia de los errores vienen de no leer bien.
    \item \textbf{Justifiquen todo:} un control sin costo-beneficio pierde puntos.
    \item \textbf{Todos deben poder explicar} cualquier parte; las preguntas cruzadas son individuales.
\end{enumerate}
\end{cajaNota}

\newpage
% ============================================================
%  SECCIÓN 7 — SIMULACRO DE PARCIAL
% ============================================================
\section{Simulacro de Parcial: caso de pr\'actica}

\begin{cajaActividad}[title={Caso: ``TurnosYa'' (practic\'a antes del Parcial)}]
\textbf{Contexto:} ``TurnosYa'' es un SaaS que gestiona turnos y fichas de pacientes para 50 cl\'inicas. Tiene 8 empleados. Funciona 100\% en la nube.\\[4pt]
\textbf{Activos principales:}
\begin{itemize}[leftmargin=1.5em]
    \item Base de datos de pacientes (incluye datos sensibles de salud).
    \item Sistema de agenda de turnos en tiempo real.
    \item Pasarela de pago de consultas.
    \item Credenciales de acceso de los m\'edicos.
\end{itemize}
\textbf{Amenazas latentes:} ransomware, ca\'ida del proveedor cloud, filtraci\'on de datos de salud, phishing a empleados.
\end{cajaActividad}

\vspace{8pt}

\noindent\textbf{Consigna del simulacro (hacelo en grupo, como si fuera el Parcial):}

\begin{enumerate}[leftmargin=2em]
    \item \textbf{Matriz de riesgos:} identific\'a al menos 3 riesgos del caso y ubicalos en la matriz P$\times$I.
    \item \textbf{BIA simplificado:} para el ``sistema de agenda de turnos'' y la ``base de datos de pacientes'', defin\'i MTD, RTO, RPO y prioridad.
    \item \textbf{3 controles priorizados:} proponelos con justificaci\'on costo-beneficio.
\end{enumerate}

\vspace{6pt}

\begin{cajaActividad}[title={Tu resoluci\'on del simulacro}]
\renewcommand{\arraystretch}{1.5}
\begin{tabularx}{\textwidth}{|X|}
\hline
\textbf{(1) Matriz de riesgos (3 riesgos con P, I y nivel):} \\[4cm]
\hline
\textbf{(2) BIA simplificado (agenda de turnos + base de datos):} \\[4cm]
\hline
\textbf{(3) 3 controles priorizados con costo-beneficio:} \\[5cm]
\hline
\end{tabularx}
\end{cajaActividad}

\newpage
% ============================================================
%  SECCIÓN 8 — GLOSARIO Y NOTAS
% ============================================================
\section{Glosario de la Clase 6}

\renewcommand{\arraystretch}{1.4}
\begin{tabularx}{\textwidth}{>{\bfseries}p{4.5cm}X}
\toprule
\textbf{T\'ermino} & \textbf{Definici\'on en una l\'inea} \\
\midrule
BIA & An\'alisis de Impacto al Negocio; mide el da\~no de una interrupci\'on. \\
MTD & M\'aximo tiempo tolerable de interrupci\'on de un proceso. \\
RTO & Objetivo de tiempo de recuperaci\'on tras un incidente. \\
RPO & P\'erdida m\'axima de datos aceptable (define frecuencia de backup). \\
Control & Medida para reducir la probabilidad o el impacto de un riesgo. \\
Control preventivo & Evita que el incidente ocurra. \\
Control detectivo & Alerta cuando el incidente est\'a ocurriendo. \\
Control correctivo & Restaura la normalidad tras el incidente. \\
Costo-beneficio & Criterio para decidir si un control vale la pena. \\
ISO/IEC 27001 & Norma de requisitos de un SGSI (certificable). \\
ISO/IEC 27002 & Cat\'alogo de buenas pr\'acticas y controles. \\
SGSI & Sistema de Gesti\'on de Seguridad de la Informaci\'on. \\
\bottomrule
\end{tabularx}

\vspace{1cm}

\section{Espacio para notas y dudas}

\begin{tcolorbox}[colback=grisClaro, colframe=gray!50, arc=2mm, height=7cm]
\textcolor{gray}{\itshape Us\'a este espacio para dudas sobre el Parcial 1, ideas de controles, o conceptos de BIA e ISO que quieras repasar.}
\end{tcolorbox}

\vfill
\begin{center}
\small\textcolor{gray}{%
Cuadernillo elaborado para la c\'atedra de Seguridad de la Informaci\'on y de
Sistemas --- Lic.\ en Gesti\'on de Negocios Digitales, UCA.\\
Prohibida su distribuci\'on fuera del \'ambito de la cursada sin autorizaci\'on del docente.}
\end{center}

\end{document}